% T17 -- Distribution et s\'ecurit\'e \'electrique, r\'egimes de terre -- banque d'exercices
%
% REFONTE 2026-09-02 : mise en compl\'ementarit\'e avec la fiche fiches/T17_Distribution_Securite.tex.
% Les trois Questions Types de la fiche sont :
%   QT1 = identification du SLT TT + choix du DDR 30 mA ;
%   QT2 = chute de tension 2024 B5-B9 (850 m, 258 A, Ru = 20 mOhm/km, Lu = 1,2 mH/km, 7,2 %) ;
%   QT3 = choix du disjoncteur NSX250N MA220 associ\'e au variateur (2021, question 17).
% Le Bloc 4 de la fiche traite en outre la s\'electivit\'e diff\'erentielle 300/30 mA (rapport >= 3)
% et l'impossibilit\'e du mode diff\'erentiel en TN-C. Les cinq exercices qui reprenaient ces
% contenus ont donc \'et\'e remplac\'es par d'autres questions des M\^EMES sujets r\'eels.
%
% Sources de la version actuelle (toutes relues dans les .md de conversion Marker) :
%   - Ex. 1 : E2_Disciplinaire_IE_2024, "\'Epreuve \'ecrite disciplinaire", questions B1 et B2.
%     Donn\'ees d'\'enonc\'e (E2_Disciplinaire_IE_2024.md l.185-201) : Un = 400 V / Vn = 230 V,
%     Is = 229 A, cos(phi_mot) = 0,85, f = 50 Hz, P0 = 1 kW avec cos(phi_0) = 1,
%     Pe = 43 kW et Qe = 25,5 kVAR. Corrig\'e officiel (Rapport_Jury_2024.md l.174-186) :
%     Pmot = 134,3 kW, Qcharge = 83 237 + 25 500 = 108,7 kVAR, Pcharge = 178,3 kW.
%     Iabs et Fp recalcul\'es (les formules du rapport sont des images omises) : 302,7 A / 0,854.
%   - Ex. 2 : E2_Sujet_2017 (EDE ENE 2, option \'energie), questions B1.1 et B1.2.
%     Puissances relev\'ees dans le corrig\'e officiel E2_Corrige_2017.md l.59 : r\'echauffeur 3 kW,
%     pompes d'hydromassage 2 x 2,25 kW, pompe de circulation 260 W. Facteurs de puissance et
%     r\'epartition sur les phases : \'enonc\'e du sujet l.147-153. Courants du corrig\'e (l.65-73) :
%     13,05 A / 11,37 A / 1,41 A -- recalcul\'es et confirm\'es.
%     (La version pr\'ec\'edente de cet exercice utilisait des puissances INVENT\'EES de 1,5 kW et
%     750 W : elles sont remplac\'ees ici par les valeurs authentiques du sujet.)
%   - Ex. 3 : E2_Disciplinaire_IE_2024, questions B3 et B4 (compensation d'\'energie r\'eactive).
%     \'Enonc\'e l.219-225 ; d\'emonstration du corrig\'e officiel Rapport_Jury_2024.md l.191-197.
%     ATTENTION : le rapport imprime "C = 0,016 F" mais sa propre formule C = Qabs/(3.omega.Vn^2)
%     donne 2,18 mF ; la valeur imprim\'ee est incoh\'erente avec la formule (artefact de la
%     conversion Marker ou coquille du rapport). On retient ici le r\'esultat de la formule.
%   - Ex. 4 : E2_Sujet_2017, question B2.3. Contexte (abonnement 18 kVA triphas\'e heures creuses,
%     16 kW de consommation r\'eguli\`ere sans d\'es\'equilibre notable) : sujet l.177.
%     Corrig\'e officiel E2_Corrige_2017.md l.105-115 : Pphase1 = 3 260 W, P = 3 x 3,26 = 9,78 kW,
%     nouvel abonnement 16 + 9,78 = 25,78 kW -> 30 kVA.
%   - Ex. 5 (docTech) : E2_Sujet_2018 (EDE ENE 2), questions 18 et 19. L'extrait DRS7 reproduit
%     dans le docTech (Imax / In / Up des PF65r, PF30r, PF30, PF15, PF8 et tableau des r\'ef\'erences
%     1P+N / 3P+N) est TEXTUELLEMENT pr\'esent dans E2_Sujet_2018.md l.888-1010 et l.1016-1038 :
%     rien n'y est reconstitu\'e. Contexte de l'installation (r\'eseau 3P+N 400 V 50 Hz, SLT TT,
%     pas de paratonnerre \`a moins de 30 m) : sujet l.243-247. Courant absorb\'e par le r\'egulateur
%     COMPACTO II (17 A/phase) : Rapport_Jury_2018.md l.2272. Corrig\'e officiel des questions 18
%     et 19 : Rapport_Jury_2018.md l.2246-2248 (PF65r 3P+N, r\'ef. 15685 ; calibre 20 A courbe C).
%     Citation jury : verbatim exact de Rapport_Jury_2018.md l.2508.
%   - Ex. 6 : E2_Sujet_2021, "\'Etude d'un syst\`eme, d'un proc\'ed\'e ou d'une organisation",
%     questions 16 et 18 (la question 17, choix du disjoncteur, est d\'ej\`a la QT3 de la fiche).
%     Corrig\'e officiel Rapport_Jury_2021.md l.637-655 : ATV630C11N4 90 kW ; In moteur = 160,4 A
%     avec eta = 0,942 et cos(phi) = 0,86 sous 400 V ; In variateur = 170 A (DT9) ; Ith = 0,94 In.
%     Citation jury : verbatim exact de Rapport_Jury_2021.md l.769 (bilan de la partie C).
%
% \'Elements SUPPRIM\'ES par cette refonte et \`a ne jamais r\'eintroduire :
%   - la citation "il est imp\'eratif de justifier le choix du disjoncteur..." attribu\'ee aux
%     rapports 2021/2024 : INTROUVABLE, fabriqu\'ee ;
%   - la configuration "300 mA en t\^ete / 30 mA par d\'epart" attribu\'ee \`a la question 19 du
%     sujet 2021 : la question 19 r\'eelle porte sur les sch\'emas de puissance et de commande ;
%   - les valeurs lin\'eiques Ru = 0,225 Ohm/km et Xu = 0,08 Ohm/km (fausses) de l'ancien
%     exercice de chute de tension.
\expandafter\ifx\csname CAPETINCLUDEMODE\endcsname\relax
\documentclass[11pt,a4paper]{article}
% preamble.tex — packages et commandes communes a toutes les fiches CAPET SII IE
\input{glyphtounicode}
\pdfgentounicode=1
\usepackage{cmap}
\usepackage[T1]{fontenc}
\usepackage[utf8]{inputenc}
\usepackage[french]{babel}
% babel[french] active a \begin{document} des raccourcis de ponctuation (:;!?)
% qui cassent la syntaxe to[...] de circuitikz (erreur "+ or - expected") ;
% on les desactive, ce qui ne change que l'espacement fin avant ces signes.
\AtBeginDocument{\shorthandoff{:;!?}}
\usepackage[a4paper, top=1.6cm, bottom=1.6cm, left=1.4cm, right=1.4cm]{geometry}
\setlength{\headheight}{14.5pt}
\usepackage{amsmath, amssymb, amsfonts}
\usepackage{siunitx}
\sisetup{locale=FR, output-decimal-marker={,}, inter-unit-product=\ensuremath{{}\cdot{}},
         group-minimum-digits=4}
% Unites non fournies par siunitx, declarees ici une fois pour toutes (fiches + exercices).
% Les gardes \ifdefined evitent tout conflit avec une redeclaration locale restee dans un
% fichier : ne pas retirer les gardes, ne pas redeclarer ces unites dans les fichiers.
\DeclareSIUnit{\jour}{j}
\ifdefined\var\else\DeclareSIUnit{\var}{var}\fi   % puissance reactive
\ifdefined\vad\else\DeclareSIUnit{\vad}{VAD}\fi   % puissance deformante
\ifdefined\tr\else\DeclareSIUnit{\tr}{tr}\fi      % tour (tr/min)
\usepackage{array, booktabs, multirow, tabularx}
\usepackage{xcolor}
\usepackage[most, breakable]{tcolorbox}
\usepackage{multicol}
\usepackage{enumitem}
\usepackage{mathtools}
\usepackage{lmodern}
% Refonte visuelle 2026 (direction "moderne coloré") : Lato pour le texte,
% maths sans-serif assorties via newtxsf (charge apres amsmath/amssymb).
\usepackage[default]{lato}
\usepackage{newtxsf}
\usepackage{microtype}
\usepackage{fancyhdr}
\usepackage{tikz}
\usetikzlibrary{positioning,arrows.meta}
\usepackage{circuitikz}
\usepackage{lastpage}

% ── Couleurs par domaine (palette modernisée 2026) ────────────────────
\definecolor{colElec}{RGB}{29, 78, 216}        % Électrotechnique — bleu vif
\definecolor{colElecLight}{RGB}{219, 234, 254} % fond bleu clair
\definecolor{colAuto}{RGB}{21, 128, 61}        % Automatique — vert
\definecolor{colMeca}{RGB}{194, 65, 12}        % Mécanique — orange
\definecolor{colTherm}{RGB}{185, 28, 28}       % Thermique — rouge
\definecolor{colQT}{RGB}{232, 244, 233}        % (hérité, plus utilisé par questiontype)
\definecolor{colJury}{RGB}{255, 248, 225}      % fond jury (crème)
\definecolor{colJuryFrame}{RGB}{202, 138, 4}   % bordure/bandeau jury (ambre)
\definecolor{colGrayLight}{RGB}{245, 245, 247} % fond gris neutre
\definecolor{colSeparator}{RGB}{180, 180, 180} % gris séparateur
\definecolor{colExo}{RGB}{8, 145, 178}          % Banque d'exercices — cyan (distinct des fiches)
% ── Encre et gris de la refonte ───────────────────────────────────────
\definecolor{colInk}{RGB}{23, 32, 42}          % texte courant (quasi-noir)
\definecolor{colSlate}{RGB}{83, 98, 118}       % gris ardoise (méta, en-têtes)
\definecolor{colRule}{RGB}{203, 213, 225}      % filets clairs
\definecolor{colJuryBg}{RGB}{255, 248, 225}    % fond crème jury
\definecolor{colPresence}{RGB}{83, 98, 118}    % bloc 5 (présence dans les sujets)
\AtBeginDocument{\color{colInk}}

% ── Style encadré exercice + corrigé (banque d'exercices) ─────────────
\tcbset{
  exoBox/.style={
    breakable,
    enhanced,
    colback=colExo!8!white,
    colframe=colExo,
    colbacktitle=colExo,
    coltitle=white,
    fonttitle=\bfseries\small,
    left=4pt, right=4pt, top=3pt, bottom=3pt,
    before skip=5pt, after skip=5pt,
  }
}

% ── Bascule enonce / corrige (banque d'exercices) ─────────────────────
% \CAPETCORRIGE est defini en ligne de commande (pdflatex "\def\CAPETCORRIGE{1}\input{...}")
% pour compiler la variante corriges d'un meme fichier source ; par defaut (non defini),
% la compilation produit la variante enonces. \CAPETSHOW{enonce}{corrige} choisit la branche.
%
% TROIS MODES (refonte 2026-09-20, manuel de revision) :
%   - aucun drapeau   -> imprime l'enonce            (variante enonces, INCHANGE)
%   - \CAPETCORRIGE   -> imprime le corrige          (variante corriges, INCHANGE)
%   - \CAPETMANUEL    -> imprime l'enonce ET met le corrige de cote pour l'annexe
%
% Le mode manuel stocke chaque corrige dans une macro (\csgdef), sans ecriture
% fichier ni re-tokenisation : le contenu (circuitikz, accents, %) traverse intact.
% L'annexe les rejoue ensuite dans l'ordre via \CAPETANNEXE.
\usepackage{etoolbox}

\newcounter{CAPETexo}          % numero continu d'exercice dans le manuel
\setcounter{CAPETexo}{0}
\newcounter{CAPETannexe}       % compteur de relecture en annexe

% Titre courant de la notion en cours (pose par \CAPETNOTION), sert d'intertitre annexe
\newcommand{\CAPETnotioncode}{}
\newcommand{\CAPETnotiontitre}{}

% \CAPETNOTION{code}{titre} : ouvre une notion dans le manuel (memorise le libelle,
% et marque le rang du premier exercice de la notion pour l'intertitre d'annexe).
\newcommand{\CAPETNOTION}[2]{%
  \renewcommand{\CAPETnotioncode}{#1}%
  \renewcommand{\CAPETnotiontitre}{#2}%
  \csgdef{CAPETnotionstart\the\numexpr\value{CAPETexo}+1\relax}{#1 -- #2}%
}

% ── \degree robuste ───────────────────────────────────────────────────
% Certains fichiers d'exercices utilisent \degree. Ni le preambule ni siunitx
% ne le definissent, et \providecommand local ne suffit pas : sous hyperref,
% \degree existe deja comme commande factice (\ERROR), donc \providecommand
% s'abstient et l'appel echoue a la compilation du manuel (erreur invisible
% en compilation isolee). On le (re)definit donc inconditionnellement ici,
% apres le chargement des packages, pour tout le corpus.
\AtBeginDocument{\providecommand{\degree}{}\renewcommand{\degree}{\ensuremath{^\circ}}}

% Libelle decoratif : les fichiers d'exercices utilisent aussi \CAPETSHOW dans
% leur \fancyhead et leur bloc de titre, pour afficher "(enonces)" ou "(corriges)".
% Ces appels-la ne sont PAS des exercices : en mode manuel ils ne doivent ni
% consommer un numero, ni alimenter l'annexe. On les reconnait a leur premier
% argument, strictement egal au libelle ci-dessous.
\newcommand{\CAPETLIBELLEENONCES}{(\'enonc\'es)}

\ifdefined\CAPETMANUEL
  % --- Mode manuel : enonce en place, corrige differe vers l'annexe ---
  \newcommand{\CAPETSHOW}[2]{%
    \def\CAPETargun{#1}%
    \ifx\CAPETargun\CAPETLIBELLEENONCES
      #1%            appel decoratif : rendu tel quel, hors numerotation
    \else
      \CAPETSHOWEXO{#1}{#2}%
    \fi
  }
  \newcommand{\CAPETSHOWEXO}[2]{%
    \stepcounter{CAPETexo}%
    \csgdef{CAPETcorrige\theCAPETexo}{#2}%
    \edef\CAPETtmpnum{\theCAPETexo}%
    #1%
    \par\vspace{2pt}%
    {\footnotesize\itshape\color{colExo}%
     \hyperref[capet:cor:\CAPETtmpnum]{Corrige de l'exercice \CAPETtmpnum\ $\rightarrow$ p.~\pageref{capet:cor:\CAPETtmpnum}}%
     \label{capet:eno:\CAPETtmpnum}}%
  }
  % \CAPETANNEXE : rejoue tous les corriges captures, dans l'ordre, avec intertitres
  \newcommand{\CAPETANNEXE}{%
    \setcounter{CAPETannexe}{0}%
    \loop\ifnum\value{CAPETannexe}<\value{CAPETexo}%
      \stepcounter{CAPETannexe}%
      \edef\CAPETcurnum{\theCAPETannexe}%
      \ifcsname CAPETnotionstart\CAPETcurnum\endcsname
        \par\medskip
        {\normalsize\bfseries\color{colExo}\csuse{CAPETnotionstart\CAPETcurnum}}%
        \par\nobreak\smallskip
      \fi
      \begin{tcolorbox}[exoBox, breakable,
        title={Corrige de l'exercice \CAPETcurnum\
               \hfill\normalfont\footnotesize
               \hyperref[capet:eno:\CAPETcurnum]{enonce p.~\pageref{capet:eno:\CAPETcurnum}}}]%
      \label{capet:cor:\CAPETcurnum}%
      \csuse{CAPETcorrige\CAPETcurnum}%
      \end{tcolorbox}%
    \repeat
  }
\else
  % --- Modes historiques, comportement strictement inchange ---
  \newcommand{\CAPETSHOW}[2]{\ifdefined\CAPETCORRIGE #2\else #1\fi}
  \newcommand{\CAPETANNEXE}{}
\fi

% ── Difficulté 1 à 3 (étoiles pleines, amssymb) ───────────────────────
\newcommand{\difficulte}[1]{%
  \ifcase#1\relax
  \or $\bigstar$\;niveau 1/3%
  \or $\bigstar\bigstar$\;niveau 2/3%
  \or $\bigstar\bigstar\bigstar$\;niveau 3/3%
  \fi
}

% ── Style encadré extrait de document technique (banque d'exercices) ──
\tcbset{
  docTech/.style={
    breakable,
    enhanced,
    colback=colGrayLight,
    colframe=colSeparator,
    fonttitle=\bfseries\small\itshape,
    coltitle=black,
    colbacktitle=colGrayLight,
    left=4pt, right=4pt, top=3pt, bottom=3pt,
    before skip=4pt, after skip=4pt,
  }
}

% ── Style tcolorbox de section principale ────────────────────────────
% Refonte 2026 : bandeau titre coloré arrondi + fond très pâle, sans cadre.
\tcbset{
  ficheSection/.style={
    breakable, enhanced jigsaw,
    frame hidden, boxrule=0pt, arc=1.8mm,
    colback=#1!4!white,
    colbacktitle=#1,
    coltitle=white,
    fonttitle=\sffamily\bfseries,
    left=7pt, right=7pt, top=5pt, bottom=5pt,
    toptitle=3pt, bottomtitle=3pt,
    before skip=10pt, after skip=8pt,
  }
}

% ── Style encadré question type ──────────────────────────────────────
% Carte blanche, bandeau sombre, liseré très clair.
\tcbset{
  questiontype/.style={
    breakable, enhanced jigsaw,
    colback=white,
    frame hidden, boxrule=0pt, arc=1.8mm,
    colbacktitle=colElec!78!black,
    coltitle=white,
    fonttitle=\sffamily\bfseries\small,
    left=7pt, right=7pt, top=4pt, bottom=4pt,
    toptitle=2.5pt, bottomtitle=2.5pt,
    before skip=8pt, after skip=8pt,
    borderline={0.6pt}{0pt}{colElec!12!white},
  }
}

% ── Style encadré jury ───────────────────────────────────────────────
% Bandeau ambre, fond crème, sans cadre.
\tcbset{
  jury/.style={
    breakable, enhanced jigsaw,
    colback=colJuryBg,
    frame hidden, boxrule=0pt, arc=1.8mm,
    colbacktitle=colJuryFrame,
    coltitle=white,
    fonttitle=\sffamily\bfseries,
    left=7pt, right=7pt, top=4pt, bottom=4pt,
    toptitle=3pt, bottomtitle=3pt,
    before skip=10pt, after skip=8pt,
  }
}

% ── En-tête / pied de page communs (refonte 2026) ─────────────────────
% #1 code fiche, #2 titre court, #3 couleur domaine
\newcommand{\ficheHeaders}[3]{%
  \fancyhead[L]{\small\sffamily{\bfseries\textcolor{#3}{#1}}\hspace{0.7em}\textcolor{colSlate}{#2}}%
  \fancyhead[R]{\small\sffamily\textcolor{colSlate}{CAPET SII IE\, --- \,Session 2026}}%
  \fancyfoot[C]{\small\sffamily\textcolor{colSlate}{\thepage\,/\,\pageref{LastPage}}}%
  \fancyhead[C]{}\fancyfoot[L]{}\fancyfoot[R]{}%
}

% ── Bloc en-tête de fiche (BLOC 0, refonte 2026) ──────────────────────
% #1 code  #2 titre  #3 couleur domaine  #4 domaine  #5 ligne présence  #6 sous-thèmes
\newcommand{\ficheEntete}[6]{%
  \begin{tcolorbox}[enhanced, breakable,
    frame hidden, boxrule=0pt, arc=2.5mm,
    interior style={left color=#3, right color=#3!75!black},
    left=9pt, right=9pt, top=7pt, bottom=7pt,
    before skip=0pt, after skip=6pt]
    {\color{white}\sffamily\bfseries\LARGE #1\quad #2}\\[4pt]
    {\color{#3!22!white}\sffamily\small \MakeUppercase{#4}\quad\textbar\quad #5}\\[1pt]
    {\color{#3!35!white}\sffamily\footnotesize Sous-thèmes : #6}
  \end{tcolorbox}%
}

% ── Mise en page ─────────────────────────────────────────────────────
\pagestyle{fancy}
\fancyhf{}
\renewcommand{\headrulewidth}{0.5pt}
\renewcommand{\footrulewidth}{0pt}
\AtBeginDocument{%
  \renewcommand{\headrule}{{\color{colRule}\hrule width\headwidth height 0.5pt}}%
}

\setlist[itemize]{leftmargin=1.2em, itemsep=1pt, topsep=2pt, parsep=0pt}
\setlist[enumerate]{leftmargin=1.2em, itemsep=3pt, topsep=2pt, parsep=0pt}

\setlength{\parskip}{2pt}
\setlength{\columnsep}{10pt}

\def\CAPETBEGINFILE{\begin{document}}
\def\CAPETENDFILE{\end{document}}
\else
\def\CAPETBEGINFILE{}
\def\CAPETENDFILE{}
\fi

\ifdefined\var\else\DeclareSIUnit{\var}{var}\fi

\fancyhead[L]{\textbf{\textcolor{colExo}{T17}} - Exercices - Distribution et s\'ecurit\'e \'electrique}
\fancyhead[R]{\small\textit{Banque d'exercices \CAPETSHOW{(\'enonc\'es)}{(corrig\'es)}}}
\fancyfoot[C]{\small\thepage\,/\,\pageref{LastPage}}

\CAPETBEGINFILE

\begin{tcolorbox}[
  enhanced, breakable,
  colback=colExo,
  colframe=colExo,
  left=6pt, right=6pt, top=5pt, bottom=5pt,
  before skip=0pt, after skip=4pt
]
  {\color{white}\bfseries\large T17 \quad Distribution et s\'ecurit\'e \'electrique, r\'egimes de terre}\\[2pt]
  {\color{white}\small Banque d'exercices -- 6 exercices \CAPETSHOW{(\'enonc\'es)}{(corrig\'es)}}
\end{tcolorbox}

% =====================================================================
% Exercice 1 -- 2024, B1 et B2
% =====================================================================
\begin{tcolorbox}[exoBox,
  title={Exercice 1 -- \difficulte{1} \hfill\normalfont\footnotesize 2024, \'epreuve \'ecrite disciplinaire (option IE), questions B1 et B2}
]
\CAPETSHOW{
L'atelier bois et menuiserie d'un lyc\'ee est aliment\'e par un r\'eseau triphas\'e \'equilibr\'e de tension compos\'ee $U_r = \SI{400}{\volt}$ (tension simple $V_n = \SI{230}{\volt}$), \`a la fr\'equence $f = \SI{50}{\hertz}$. Il regroupe trois groupes de r\'ecepteurs :

\begin{itemize}[leftmargin=1.5em, itemsep=1pt]
  \item le \textbf{moto-ventilateur} de l'extracteur d'air, machine asynchrone triphas\'ee de courant nominal $I_S = \SI{229}{\ampere}$ sous $U_n = \SI{400}{\volt}$, de facteur de puissance $\cos\varphi_{mot} = 0{,}85$ ;
  \item les \textbf{ordinateurs et l'automate de s\'ecurit\'e}, qui consomment $P_0 = \SI{1}{\kilo\watt}$ avec $\cos\varphi_0 = 1$ ;
  \item l'\textbf{\'eclairage et les \'equipements} de l'atelier (groupe de per\c{c}age, brosses, cha\^ines et scies), qui consomment $P_e = \SI{43}{\kilo\watt}$ et $Q_e = \SI{25,5}{\kilo\var}$.
\end{itemize}

\begin{center}
\begin{tikzpicture}[>=Stealth, scale=0.9, every node/.style={scale=0.9}]
  \foreach \y in {2.4,2.1,1.8} { \draw[thick] (0,\y) -- (8.1,\y); }
  \node[left, align=center, font=\scriptsize] at (0,2.1) {R\'eseau\\triphas\'e\\\SI{400}{\volt}};
  \node[draw, minimum width=1.9cm, minimum height=0.9cm, align=center, font=\scriptsize] (m) at (1.9,0.5) {MAS\\extracteur\\$I_S{=}\SI{229}{\ampere}$};
  \node[draw, minimum width=1.9cm, minimum height=0.9cm, align=center, font=\scriptsize] (o) at (4.4,0.5) {Ordinateurs\\+ automate\\$P_0{=}\SI{1}{\kilo\watt}$};
  \node[draw, minimum width=2.1cm, minimum height=0.9cm, align=center, font=\scriptsize] (e) at (7.0,0.5) {\'Eclairage\\+ \'equipements\\$P_e{=}\SI{43}{\kilo\watt}$};
  \foreach \n in {m,o,e} { \draw (\n.north) -- (\n.north |- 0,1.8); }
\end{tikzpicture}
\end{center}

\begin{enumerate}[leftmargin=1.5em]
  \item Calculer la puissance active $P_{mot}$ et la puissance r\'eactive $Q_{mot}$ absorb\'ees par le moto-ventilateur.
  \item En appliquant le th\'eor\`eme de Boucherot, en d\'eduire la puissance active totale $P_{charge}$, la puissance r\'eactive totale $Q_{charge}$ et la puissance apparente totale $S_{charge}$ de l'atelier.
  \item En d\'eduire le courant total absorb\'e $I_{abs}$ par l'installation, ainsi que son facteur de puissance $F_p$.
  \item Le gestionnaire d'\'energie impose un facteur de puissance compris entre $0{,}92$ et $1$. Conclure.
\end{enumerate}

}{
\textbf{Corrig\'e.}

\begin{enumerate}[leftmargin=1.5em]
  \item La machine \'etant triphas\'ee \'equilibr\'ee et aliment\'ee sous la tension simple $V_n$ :
  \[ P_{mot} = 3\,V_n\,I_S\cos\varphi_{mot} = 3\times 230\times 229\times 0{,}85 \approx \SI{134,3}{\kilo\watt} \]
  \[ Q_{mot} = P_{mot}\tan\varphi_{mot} = 134\,308\times\tan(\arccos 0{,}85) \approx \SI[group-minimum-digits=4]{83237}{\var} \approx \SI{83,2}{\kilo\var} \]
  \item Le th\'eor\`eme de Boucherot autorise la somme s\'epar\'ee des puissances actives et des puissances r\'eactives (jamais des puissances apparentes). Les ordinateurs et l'automate ayant $\cos\varphi_0 = 1$, ils n'apportent aucun r\'eactif ($Q_0 = 0$) :
  \[ P_{charge} = P_{mot} + P_0 + P_e = 134{,}3 + 1 + 43 = \boldsymbol{\SI{178,3}{\kilo\watt}} \]
  \[ Q_{charge} = Q_{mot} + Q_0 + Q_e = 83{,}2 + 0 + 25{,}5 = \boldsymbol{\SI{108,7}{\kilo\var}} \]
  \[ S_{charge} = \sqrt{P_{charge}^2 + Q_{charge}^2} = \sqrt{178{,}3^2 + 108{,}7^2} \approx \SI{208,8}{\kilo\volt\ampere} \]
  \item L'installation \'etant triphas\'ee, $S_{charge} = 3\,V_n\,I_{abs} = \sqrt{3}\,U_r\,I_{abs}$, d'o\`u :
  \[ I_{abs} = \frac{S_{charge}}{3\,V_n} = \frac{208\,800}{3\times 230} \approx \SI{303}{\ampere} \]
  \[ F_p = \frac{P_{charge}}{S_{charge}} = \frac{178{,}3}{208{,}8} \approx 0{,}854 \]
  \item \textbf{Conclusion.} Le facteur de puissance mesur\'e, $F_p \approx 0{,}854$, est \textbf{inf\'erieur \`a la borne basse de $0{,}92$} impos\'ee par le gestionnaire d'\'energie : l'exigence contractuelle n'est pas satisfaite en l'\'etat. L'installation transporte $\SI{108,7}{\kilo\var}$ de r\'eactif qui ne produit aucun travail utile mais surcharge la ligne (303 A au lieu des 258 A qu'appellerait la m\^eme puissance active \`a $F_p = 1$). Une \textbf{batterie de condensateurs de compensation} doit donc \^etre ajout\'ee en parall\`ele de l'installation : c'est l'objet de l'exercice 3.
\end{enumerate}
}
\end{tcolorbox}

% =====================================================================
% Exercice 2 -- 2017, B1.1 et B1.2
% =====================================================================
\begin{tcolorbox}[exoBox,
  title={Exercice 2 -- \difficulte{1} \hfill\normalfont\footnotesize 2017, EDE ENE 2 (option \'energie), questions B1.1 et B1.2}
]
\CAPETSHOW{
L'armoire \'electrique d'un spa domestique r\'epartit ses r\'ecepteurs \textbf{monophas\'es} (tension simple $U = \SI{230}{\volt}$) sur les trois phases du r\'eseau triphas\'e. La documentation constructeur du spa donne les puissances suivantes : r\'echauffeur \SI{3}{\kilo\watt} ; pompes d'hydromassage $2\times\SI{2,25}{\kilo\watt}$ ; pompe de circulation \SI{260}{\watt}.

Le sujet pr\'ecise la r\'epartition et les facteurs de puissance :

\begin{itemize}[leftmargin=1.5em, itemsep=1pt]
  \item le \textbf{r\'echauffeur} (purement r\'esistif, $\cos\varphi = 1$) est aliment\'e par la \textbf{phase 1} ;
  \item la \textbf{pompe de circulation} ($\cos\varphi = 0{,}80$), utilis\'ee pour la filtration et le r\'echauffeur, est \'egalement aliment\'ee par la \textbf{phase 1} ;
  \item les \textbf{pompes d'hydromassage 1 et 2} ($\cos\varphi = 0{,}86$) sont aliment\'ees respectivement par les \textbf{phases 2 et 3}, et fonctionnent le plus souvent simultan\'ement.
\end{itemize}

\begin{center}
\begin{tikzpicture}[>=Stealth, scale=0.85, every node/.style={scale=0.85}]
  \foreach \p/\y in {1/2.6,2/1.4,3/0.2} {
    \draw[thick] (0,\y) -- (7.4,\y) node[right,font=\scriptsize]{Phase \p};
  }
  \node[draw,minimum width=2.1cm,minimum height=0.6cm,align=center,font=\scriptsize] (rc) at (1.6,3.4) {R\'echauffeur\\\SI{3}{\kilo\watt}, $\cos\varphi{=}1$};
  \draw (1.6,3.1) -- (1.6,2.6);
  \node[draw,minimum width=2.4cm,minimum height=0.6cm,align=center,font=\scriptsize] (pc) at (4.8,3.4) {Pompe circulation\\\SI{260}{\watt}, $\cos\varphi{=}0{,}80$};
  \draw (4.8,3.1) -- (4.8,2.6);
  \node[draw,minimum width=2.4cm,minimum height=0.6cm,align=center,font=\scriptsize] (p1) at (3.0,2.0) {Pompe hydro 1\\\SI{2,25}{\kilo\watt}, $\cos\varphi{=}0{,}86$};
  \draw (3.0,1.7) -- (3.0,1.4);
  \node[draw,minimum width=2.4cm,minimum height=0.6cm,align=center,font=\scriptsize] (p2) at (3.0,0.8) {Pompe hydro 2\\\SI{2,25}{\kilo\watt}, $\cos\varphi{=}0{,}86$};
  \draw (3.0,0.5) -- (3.0,0.2);
\end{tikzpicture}
\end{center}

\begin{enumerate}[leftmargin=1.5em]
  \item Rappeler l'expression du courant absorb\'e par un r\'ecepteur \textbf{monophas\'e} en fonction de sa puissance active, de sa tension d'alimentation et de son facteur de puissance.
  \item Calculer le courant consomm\'e par chacun des quatre r\'ecepteurs.
  \item Le r\'echauffeur, le plus gros consommateur, sera prot\'eg\'e par un appareil de commande et de protection int\'egral dont les calibres disponibles sont : $1{,}25$ \`a \SI{5}{\ampere} ; $3$ \`a \SI{12}{\ampere} ; $4{,}5$ \`a \SI{18}{\ampere}. Indiquer le calibre \`a retenir et conclure sur l'ad\'equation de ce choix.
\end{enumerate}

}{
\textbf{Corrig\'e.}

\begin{enumerate}[leftmargin=1.5em]
  \item Pour un r\'ecepteur monophas\'e, la puissance active vaut $P = U\,I\cos\varphi$, d'o\`u :
  \[ I = \frac{P}{U\cos\varphi} \]
  Le facteur $\cos\varphi$ est bien au \textbf{d\'enominateur} : plus il est faible, plus le courant appel\'e est \'elev\'e \`a puissance active \'egale.
  \item \[ I_{\text{r\'echauffeur}} = \frac{3\,000}{230\times 1} \approx \SI{13,05}{\ampere} \qquad I_{\text{pompe circ.}} = \frac{260}{230\times 0{,}80} \approx \SI{1,41}{\ampere} \]
  \[ I_{\text{pompe hydro 1}} = I_{\text{pompe hydro 2}} = \frac{2\,250}{230\times 0{,}86} \approx \SI{11,37}{\ampere} \]
  On remarque que la pompe de circulation, malgr\'e son mauvais facteur de puissance ($0{,}80$), reste le r\'ecepteur le moins consommateur : sa faible puissance active domine l'effet du $\cos\varphi$.
  \item Le courant du r\'echauffeur vaut $\SI{13,05}{\ampere}$. Il faut donc un calibre dont la \textbf{plage de r\'eglage encadre} cette valeur :
  \begin{itemize}
    \item $1{,}25$ \`a \SI{5}{\ampere} : trop faible, $13{,}05 > 5$ -- d\'eclenchement permanent ;
    \item $3$ \`a \SI{12}{\ampere} : trop faible \'egalement, $13{,}05 > 12$ ;
    \item $4{,}5$ \`a \SI{18}{\ampere} : \textbf{convient}, $4{,}5 < 13{,}05 < 18$.
  \end{itemize}

  \smallskip
  \textbf{Conclusion.} On retient le calibre \textbf{$4{,}5$ \`a \SI{18}{\ampere}}, seul \`a encadrer le courant d'emploi de \SI{13,05}{\ampere} : le r\'eglage de la protection thermique pourra \^etre positionn\'e sur la valeur exacte du courant d'emploi, condition impos\'ee par la NF C 15-100 pour la protection contre les surcharges ($I_B \leq I_n \leq I_z$). Le calibre imm\'ediatement inf\'erieur ($3$ \`a \SI{12}{\ampere}) aurait provoqu\'e des d\'eclenchements intempestifs en fonctionnement nominal.
\end{enumerate}
}
\end{tcolorbox}

% =====================================================================
% Exercice 3 -- 2024, B3 et B4
% =====================================================================
\begin{tcolorbox}[exoBox,
  title={Exercice 3 -- \difficulte{2} \hfill\normalfont\footnotesize 2024, \'epreuve \'ecrite disciplinaire (option IE), questions B3 et B4}
]
\CAPETSHOW{
On reprend l'atelier de l'exercice 1, dont le bilan de puissances a conduit \`a $P_{charge} = \SI{178,3}{\kilo\watt}$ et $Q_{charge} = \SI{108,7}{\kilo\var}$ sous une tension simple $V_n = \SI{230}{\volt}$ \`a $f = \SI{50}{\hertz}$.

Pour respecter l'exigence du gestionnaire d'\'energie, une \textbf{batterie de condensateurs} est ajout\'ee \textbf{en parall\`ele} de l'installation : un condensateur de capacit\'e $C$ par phase, coupl\'es en \textbf{\'etoile}. Le dimensionnement est r\'ealis\'e de mani\`ere \`a ramener le facteur de puissance de l'ensemble \`a la valeur \textbf{unitaire} au point de fonctionnement nominal.

\begin{center}
\begin{circuitikz}[european, scale=0.85, every node/.style={scale=0.85}]
  \draw (0,3.0) node[left,font=\scriptsize]{Ph 1} -- (5.6,3.0);
  \draw (0,2.6) node[left,font=\scriptsize]{Ph 2} -- (4.4,2.6);
  \draw (0,2.2) node[left,font=\scriptsize]{Ph 3} -- (3.2,2.2);
  % batterie de condensateurs en etoile
  \draw (3.2,2.2) to[C, l_=$C$] (3.2,0.6);
  \draw (4.4,2.6) to[C, l_=$C$] (4.4,0.6);
  \draw (5.6,3.0) to[C, l=$C$] (5.6,0.6);
  \draw (3.2,0.6) -- (5.6,0.6);
  \draw (4.4,0.6) -- (4.4,0.35);
  \node[font=\scriptsize, below] at (4.4,0.3) {point neutre de la batterie};
  % charge
  \node[draw, minimum width=2.0cm, minimum height=0.9cm, align=center, font=\scriptsize] (ch) at (7.6,1.7) {Atelier\\$P_{charge}$\\$Q_{charge}$};
  \draw (5.6,3.0) -- (7.6,3.0) -- (ch.north);
  \fill (3.2,2.2) circle (1.4pt); \fill (4.4,2.6) circle (1.4pt); \fill (5.6,3.0) circle (1.4pt);
\end{circuitikz}
\end{center}

\begin{enumerate}[leftmargin=1.5em]
  \item \'Ecrire la condition portant sur la puissance r\'eactive totale (charge + batterie) pour que le facteur de puissance de l'ensemble soit unitaire.
  \item Rappeler l'expression de la puissance r\'eactive $Q_C$ fournie par une batterie triphas\'ee de trois condensateurs $C$ coupl\'es en \'etoile sous la tension simple $V_n$. En d\'eduire l'expression litt\'erale de $C$ en fonction de $Q_{charge}$, $\omega$ et $V_n$.
  \item Calculer la valeur num\'erique de $C$ et conclure sur la faisabilit\'e technique de cette compensation.
\end{enumerate}

}{
\textbf{Corrig\'e.}

\begin{enumerate}[leftmargin=1.5em]
  \item Le facteur de puissance de l'ensemble vaut $F_p' = P_{totale}/S_{totale}$. Les condensateurs \'etant suppos\'es parfaits, ils ne consomment aucune puissance active : $P_{totale} = P_{charge}$. La condition $F_p' = 1$ \'equivaut donc \`a $S_{totale} = P_{totale}$, c'est-\`a-dire :
  \[ Q_{totale} = Q_C + Q_{charge} = 0 \qquad\Longleftrightarrow\qquad \boldsymbol{Q_C = -\,Q_{charge}} \]
  Les condensateurs doivent \textbf{fournir} exactement la puissance r\'eactive que la charge \textbf{absorbe} : la circulation de r\'eactif est alors confin\'ee \`a la boucle locale charge/batterie et dispara\^it de la ligne d'alimentation.
  \item Chaque condensateur, coupl\'e en \'etoile, est soumis \`a la tension simple $V_n$ et absorbe une puissance r\'eactive $-C\,\omega\,V_n^2$ (n\'egative : un condensateur fournit du r\'eactif). Pour les trois phases :
  \[ Q_C = -\,3\,C\,\omega\,V_n^2 \]
  En reportant dans la condition pr\'ec\'edente, $-3\,C\,\omega\,V_n^2 = -Q_{charge}$, d'o\`u :
  \[ \boxed{\;C = \frac{Q_{charge}}{3\,\omega\,V_n^2}\;} \]
  \item Avec $\omega = 2\pi f = 2\pi\times 50 \approx \SI{314}{\radian\per\second}$ :
  \[ C = \frac{108{,}7\times 10^3}{3\times 2\pi\times 50\times 230^2} = \frac{108\,700}{3\times 314\times 52\,900} \approx \SI{2,18e-3}{\farad} \approx \SI{2,2}{\milli\farad} \]

  \smallskip
  \textbf{Conclusion.} Il faut trois condensateurs d'environ \SI{2,2}{\milli\farad} (soit \SI{2200}{\micro\farad}) tenant la tension simple de \SI{230}{\volt}, ce qui correspond \`a des mod\`eles de puissance courants du commerce (une batterie de compensation d'environ \SI{36}{\kilo\var} par phase) : la compensation est donc \textbf{techniquement r\'ealisable}. Une fois install\'ee, l'atelier est vu par le r\'eseau comme une charge purement active de \SI{178,3}{\kilo\watt} \`a $F_p = 1$, ce qui satisfait l'exigence contractuelle du gestionnaire d'\'energie ($0{,}92 \leq F_p \leq 1$) et ram\`ene le courant de ligne de \SI{303}{\ampere} \`a \SI{258}{\ampere}, soit un all\`egement de pr\`es de \SI{15}{\percent} sur le dimensionnement des c\^ables.

  \smallskip
  \textbf{Remarque.} La compensation ne r\'eduit pas la puissance active consomm\'ee, donc pas la facture d'\'energie au sens strict : elle supprime la p\'enalit\'e sur le r\'eactif et lib\`ere de la capacit\'e de transport sur l'installation existante.
\end{enumerate}
}
\end{tcolorbox}

% =====================================================================
% Exercice 4 -- 2017, B2.3
% =====================================================================
\begin{tcolorbox}[exoBox,
  title={Exercice 4 -- \difficulte{2} \hfill\normalfont\footnotesize 2017, EDE ENE 2 (option \'energie), question B2.3}
]
\CAPETSHOW{
On reprend l'installation du spa de l'exercice 2. Le propri\'etaire dispose d\'ej\`a d'un \textbf{abonnement de \SI{18}{\kilo\volt\ampere} en triphas\'e} avec heures creuses ; ses derni\`eres factures font appara\^itre des consommations r\'eguli\`eres de \textbf{\SI{16}{\kilo\watt} au maximum, sans d\'es\'equilibre notable}.

L'installation du spa vient s'ajouter \`a cette consommation existante. Le cas de fonctionnement le plus contraignant est celui o\`u le r\'echauffeur fonctionne : le r\'echauffeur et la pompe de circulation sont alors tous deux aliment\'es par la \textbf{phase 1}, tandis que les phases 2 et 3 ne re\c{c}oivent rien (les pompes d'hydromassage sont \`a l'arr\^et d\`es que le r\'echauffeur fonctionne).

On rappelle qu'en triphas\'e la puissance souscrite est \textbf{r\'epartie \`a parts \'egales entre les trois phases} : le disjoncteur de branchement d\'eclenche d\`es qu'\textbf{une seule} phase d\'epasse le tiers de la puissance souscrite, quel que soit le niveau de charge des deux autres.

Les puissances souscrites normalis\'ees disponibles en triphas\'e sont : $3$, $6$, $9$, $12$, $15$, $18$, $24$ et \SI{30}{\kilo\volt\ampere}.

\begin{enumerate}[leftmargin=1.5em]
  \item Calculer la puissance suppl\'ementaire $P_{phase\,1}$ appel\'ee sur la phase 1 par le spa dans ce cas de fonctionnement.
  \item En tenant compte de la r\`egle de r\'epartition rappel\'ee ci-dessus, en d\'eduire la puissance suppl\'ementaire $P$ qu'il faut souscrire sur l'abonnement \textbf{triphas\'e}. Expliquer pourquoi ce n'est pas simplement $P_{phase\,1}$.
  \item D\'eterminer la puissance totale n\'ecessaire, puis le nouvel abonnement \`a souscrire. Conclure sur la pertinence \'economique du choix du triphas\'e dans cette situation.
\end{enumerate}

}{
\textbf{Corrig\'e.}

\begin{enumerate}[leftmargin=1.5em]
  \item C'est lors du fonctionnement du r\'echauffeur que la puissance appel\'ee est la plus \'elev\'ee. Le r\'echauffeur et la pompe de circulation \'etant tous deux sur la phase 1 :
  \[ P_{phase\,1} = P_{\text{r\'echauffeur}} + P_{\text{circulation}} = 3\,000 + 260 = \SI[group-minimum-digits=4]{3260}{\watt} = \SI{3,26}{\kilo\watt} \]
  \item Le disjoncteur de branchement triphas\'e surveille \textbf{chaque phase ind\'ependamment}, et la puissance souscrite est divis\'ee par trois pour d\'efinir le seuil de chaque phase. Pour que la phase 1 puisse encaisser \SI{3,26}{\kilo\watt} de plus, il faut donc augmenter la puissance souscrite de \textbf{trois fois} cette valeur :
  \[ P = 3\times P_{phase\,1} = 3\times 3{,}26 = \SI{9,78}{\kilo\watt} \]
  Ce n'est pas simplement $P_{phase\,1}$ car la souscription est un contrat \textbf{global} : on ne peut pas acheter de la puissance sur une seule phase. Les \SI{6,52}{\kilo\watt} suppl\'ementaires souscrits sur les phases 2 et 3 seront in\'evitablement pay\'es sans jamais \^etre utilis\'es -- c'est le co\^ut direct du \textbf{d\'es\'equilibre} de l'installation.
  \item Le nouvel abonnement doit couvrir la consommation existante augment\'ee de ce suppl\'ement :
  \[ P_{totale} = 16 + 9{,}78 = \SI{25,78}{\kilo\watt} \]
  La valeur normalis\'ee imm\'ediatement sup\'erieure \'etant \SI{30}{\kilo\volt\ampere} (celle de \SI{24}{\kilo\volt\ampere} \'etant insuffisante), il faut souscrire un \textbf{nouvel abonnement de \SI{30}{\kilo\volt\ampere}}.

  \smallskip
  \textbf{Conclusion.} Le d\'es\'equilibre de l'installation force \`a passer de \SI{18}{\kilo\volt\ampere} \`a \SI{30}{\kilo\volt\ampere}, soit deux crans normalis\'es, alors que la puissance r\'eellement ajout\'ee n'est que de \SI{3,26}{\kilo\watt} : environ \SI{6,5}{\kilo\watt} de puissance souscrite seront pay\'es sans jamais \^etre consomm\'es. C'est un crit\`ere de conception d\'eterminant : \textbf{r\'e\'equilibrer la r\'epartition} des r\'ecepteurs sur les trois phases (par exemple en d\'epla\c{c}ant la pompe de circulation sur la phase 2) aurait limit\'e le suppl\'ement \`a $3\times 3 = \SI{9}{\kilo\watt}$, et une conception plus fine du d\'ecoupage aurait pu \'eviter le changement d'abonnement.
\end{enumerate}
}
\end{tcolorbox}

% =====================================================================
% Exercice 5 -- docTech -- 2018, questions 18 et 19
% =====================================================================
\begin{tcolorbox}[exoBox,
  title={Exercice 5 -- \difficulte{3} \hfill\normalfont\footnotesize 2018, EDE ENE 2, questions 18 et 19 (document ressource DRS7)}
]
\CAPETSHOW{
L'\'eclairage public de la rue principale d'une commune est pilot\'e par un r\'egulateur-variateur COMPACTO~II\textsuperscript{\textregistered}, install\'e dans un coffret en bord de voirie. Le dossier technique pr\'ecise :

\begin{itemize}[leftmargin=1.5em, itemsep=1pt]
  \item r\'eseau \'electrique d'alimentation triphas\'e \textbf{3 P + N}, \SI{400}{\volt}, \SI{50}{\hertz} ;
  \item sch\'ema de liaison \`a la terre de l'installation de type \textbf{TT} ;
  \item pas de paratonnerre \`a moins de \SI{30}{\meter} du coffret d'installation.
\end{itemize}

Le constructeur pr\'econise l'installation de parafoudres de \textbf{type 2} en amont et en aval du r\'egulateur, assurant la protection \`a la fois contre les surtensions en \textbf{mode commun (MC)} et en \textbf{mode diff\'erentiel (MD)}. Le \textbf{courant maximum de d\'echarge est estim\'e \`a \SI{15}{\kilo\ampere}}. Le r\'egulateur absorbe \SI{17}{\ampere} par phase.

\begin{tcolorbox}[docTech, title={Document technique DRS7 -- parafoudres de type 2, gamme PF (d'apr\`es Merlin Gerin)}]
\footnotesize
\textbf{Caract\'eristiques par mod\`ele} (onde de courant normalis\'ee $8/20~\upmu\text{s}$) :

\begin{center}
\begin{tabular}{@{}lccc@{}}
\toprule
Mod\`ele & Protection & $I_{max}$ ($8/20~\upmu$s) & $U_p$ \\
\midrule
PF 65r & mode commun & \SI{65}{\kilo\ampere} & \SI{2}{\kilo\volt} \\
PF 30r / PF 30 & mode commun & \SI{30}{\kilo\ampere} & \SI{1,8}{\kilo\volt} \\
PF 15 & mode commun & \SI{15}{\kilo\ampere} & \SI{1,8}{\kilo\volt} \\
       & mode diff\'erentiel & \SI{8}{\kilo\ampere} & \SI{1}{\kilo\volt} \\
PF 8  & mode commun & \SI{8}{\kilo\ampere} & \SI{1,5}{\kilo\volt} \\
       & mode diff\'erentiel & \SI{8}{\kilo\ampere} & \SI{1}{\kilo\volt} \\
\bottomrule
\end{tabular}
\end{center}

\emph{Le PF65r int\`egre une protection en mode diff\'erentiel ainsi qu'un contact de signalisation \`a distance ; il est conseill\'e pour un niveau de risque tr\`es \'elev\'e. Le PF8 assure une protection fine des r\'ecepteurs et se place en cascade avec les parafoudres de t\^ete.}

\smallskip
\textbf{R\'ef\'erences PF} ($U_n = \SI{400}{\volt}$) :

\begin{center}
\begin{tabular}{@{}llccc@{}}
\toprule
Type & Mod\`ele & $U_c$ (V) & $I_n$ MC/MD (kA) & R\'ef. \\
\midrule
1P + N & PF 65r & 440 & 20 & 15\,684 \\
       & PF 30r & 440 & 10 & 15\,689 \\
       & PF 30  & 440 & 10 & 15\,687 \\
       & PF 15  & 440/250 & 5 & 15\,692 \\
       & PF 8   & 440/250 & 2 & 15\,695 \\
\midrule
3P + N & PF 65r & 440 & 20 & 15\,685 \\
       & PF 30r & 440 & 10 & 15\,690 \\
       & PF 30  & 440 & 10 & 15\,688 \\
       & PF 15  & 440/250 & 5 & 15\,693 \\
       & PF 8   & 440/250 & 2 & 15\,696 \\
\bottomrule
\end{tabular}
\end{center}

\smallskip
\textbf{D\'econnexion des parafoudres.} \emph{Un des param\`etres du parafoudre est la valeur du courant maximum $I_{max}$ (onde $8/20~\upmu$s) qu'il peut supporter sans d\'egradation. Si cette valeur est d\'epass\'ee, le parafoudre sera d\'etruit et se mettra en court-circuit d\'efinitivement. Le courant de d\'efaut devra \^etre \'elimin\'e par un d\'econnecteur externe plac\'e en amont. Seul un disjoncteur pourra assurer la protection compl\`ete exig\'ee : ne pas d\'eclencher \`a 20 chocs \`a $I_n$, pouvoir d\'eclencher \`a $I_{max}$ sans se d\'etruire, et assurer la d\'econnexion du parafoudre en cas de court-circuit de ce dernier.}
\end{tcolorbox}

\begin{enumerate}[leftmargin=1.5em]
  \item \textbf{Relever} dans le document technique, pour chacun des cinq mod\`eles de la gamme PF, la valeur du courant maximum de d\'echarge $I_{max}$ en mode commun, et \textbf{indiquer} lesquels de ces mod\`eles offrent en outre une protection en mode diff\'erentiel.
  \item En d\'eduire, en justifiant par les deux crit\`eres impos\'es par le cahier des charges (courant maximum de d\'echarge estim\'e et modes de protection exig\'es), le \textbf{mod\`ele} de parafoudre \`a retenir. Attention au pi\`ege de d\'esignation de la gamme.
  \item \textbf{D\'eterminer} la r\'ef\'erence compl\`ete du parafoudre \`a commander, en tenant compte de la nature du r\'eseau d'alimentation.
  \item Le disjoncteur de d\'econnexion associ\'e doit avoir le \textbf{m\^eme calibre} que le disjoncteur de t\^ete, de mani\`ere \`a ne pas cr\'eer de s\'electivit\'e entre eux. Les calibres normalis\'es disponibles sont $10$, $16$, $20$, $25$ et \SI{32}{\ampere}, en courbes B, C ou D. \textbf{D\'eterminer} le calibre et le type de courbe de d\'eclenchement \`a retenir, puis conclure.
\end{enumerate}

}{
\textbf{Corrig\'e.}

\begin{enumerate}[leftmargin=1.5em]
  \item Lecture directe du tableau des caract\'eristiques :
  \begin{center}
  \begin{tabular}{@{}lcc@{}}
  \toprule
  Mod\`ele & $I_{max}$ mode commun & Protection mode diff\'erentiel \\
  \midrule
  PF 65r & \SI{65}{\kilo\ampere} & oui \\
  PF 30r / PF 30 & \SI{30}{\kilo\ampere} & non \\
  PF 15 & \SI{15}{\kilo\ampere} & oui (mais $I_{max} = \SI{8}{\kilo\ampere}$ en MD) \\
  PF 8 & \SI{8}{\kilo\ampere} & oui (mais $I_{max} = \SI{8}{\kilo\ampere}$ en MD) \\
  \bottomrule
  \end{tabular}
  \end{center}
  Seuls les \textbf{PF 65r}, \textbf{PF 15} et \textbf{PF 8} offrent les deux modes de protection ; les PF 30r et PF 30 ne prot\`egent qu'en mode commun.
  \item Deux crit\`eres doivent \^etre satisfaits \textbf{simultan\'ement} :
  \begin{itemize}
    \item \textbf{Courant maximum de d\'echarge} : le parafoudre doit supporter $I_m = \SI{15}{\kilo\ampere}$ \emph{sans d\'egradation}. Le PF 8 ($I_{max} = \SI{8}{\kilo\ampere}$) est \'elimin\'e. Le PF 15 affiche bien $I_{max} = \SI{15}{\kilo\ampere}$ en mode commun, \textbf{mais seulement \SI{8}{\kilo\ampere} en mode diff\'erentiel} : il ne tient donc pas les \SI{15}{\kilo\ampere} sur les deux modes.
    \item \textbf{Modes de protection} : MC \emph{et} MD sont exig\'es par le constructeur du r\'egulateur. Les PF 30r et PF 30 sont \'elimin\'es (mode commun seul).
  \end{itemize}
  Le seul mod\`ele satisfaisant les deux crit\`eres est le \textbf{PF 65r} ($I_{max} = \SI{65}{\kilo\ampere} \geq \SI{15}{\kilo\ampere}$, protection MC et MD, signalisation \`a distance).

  \smallskip
  \emph{Pi\`ege de d\'esignation} : le chiffre du nom commercial (PF \textbf{15}) d\'esigne l'ordre de grandeur du courant maximum de d\'echarge du mod\`ele, \textbf{pas} le courant de d\'echarge de l'installation. Choisir le PF 15 « parce que $I_m = \SI{15}{\kilo\ampere}$ » revient \`a dimensionner le composant exactement \`a sa limite de destruction, sans aucune marge, et sans satisfaire le mode diff\'erentiel.
  \item Le r\'eseau \'etant \textbf{triphas\'e 3 P + N}, il faut la version \textbf{3P + N} du PF 65r (et non la version 1P + N, monophas\'ee). D'apr\`es le tableau des r\'ef\'erences :
  \[ \text{PF 65r, 3P + N} \;\longrightarrow\; \text{r\'ef\'erence } \boldsymbol{15\,685} \]
  \item Le disjoncteur de d\'econnexion doit \^etre adapt\'e au parafoudre et, pour \'eviter toute s\'electivit\'e avec le disjoncteur de t\^ete, porter le \textbf{m\^eme calibre} que celui-ci. Le r\'egulateur COMPACTO~II\textsuperscript{\textregistered} absorbant $I = \SI{17}{\ampere}$ par phase, le calibre normalis\'e imm\'ediatement sup\'erieur est \textbf{\SI{20}{\ampere}} (le calibre \SI{16}{\ampere} serait inf\'erieur au courant d'emploi).

  Le choix de la courbe se fait sur la nature de la charge : la charge est un r\'egulateur d'\'eclairage, sans moteur donc \textbf{sans courant de d\'emarrage} et \textbf{sans d\'epart long} (pas de forte pointe d'enclenchement \`a tol\'erer). La \textbf{courbe C} (seuil magn\'etique entre $5$ et $10\,I_n$) est donc adapt\'ee ; la courbe D, r\'eserv\'ee aux charges \`a fort courant d'appel, serait inutilement permissive.

  \smallskip
  \textbf{Conclusion.} L'ensemble retenu -- \textbf{parafoudre PF 65r 3P + N, r\'ef. 15\,685}, associ\'e \`a un \textbf{disjoncteur de d\'econnexion \SI{20}{\ampere} courbe C} -- satisfait les trois exigences du cahier des charges : il supporte le courant de d\'echarge estim\'e de \SI{15}{\kilo\ampere} avec une large marge, prot\`ege le r\'egulateur en mode commun \emph{et} en mode diff\'erentiel comme le pr\'econise le constructeur, et permet la d\'econnexion du parafoudre en fin de vie sans provoquer la coupure s\'elective de l'installation compl\`ete.

  \smallskip
  \textit{Point de vigilance jury (rapport de jury CAPET 2018, commentaire des questions 16 \`a 19) : \guillemotleft\ C'est ici le sch\'ema de liaison \`a la terre de type TT qui conditionnait et justifiait le choix des diff\'erents \'el\'ements de protection. Son importance n'a pas \'et\'e relev\'ee par un grand nombre de candidats. [\dots] Certains candidats ont \'egalement confondu la r\'ef\'erence du parafoudre PF15 avec le courant de d\'echarge de 15 kA. Le mod\`ele est rarement donn\'e avec la r\'ef\'erence en triphas\'e et certains candidats font un choix de parafoudre en monophas\'e. \guillemotright}
\end{enumerate}
}
\end{tcolorbox}

% =====================================================================
% Exercice 6 -- 2021, questions 16 et 18
% =====================================================================
\begin{tcolorbox}[exoBox,
  title={Exercice 6 -- \difficulte{3} \hfill\normalfont\footnotesize 2021, \'Etude d'un syst\`eme, d'un proc\'ed\'e ou d'une organisation, questions 16 et 18}
]
\CAPETSHOW{
Le groupe motopompe d'une micro-STEP solaire est entra\^in\'e par un moteur asynchrone triphas\'e de \textbf{puissance utile $P_u = \SI{90}{\kilo\watt}$}, aliment\'e sous une tension compos\'ee $U = \SI{400}{\volt}$. Sa plaque signal\'etique donne un rendement $\eta = 0{,}942$ et un facteur de puissance $\cos\varphi = 0{,}86$.

Ce moteur est pilot\'e par un \textbf{variateur de vitesse Schneider ATV630}. Le document technique DT9 du sujet indique que, pour un moteur de \SI{90}{\kilo\watt} sous \SI{400}{\volt} avec une capacit\'e de surcharge de \SI{150}{\percent} du courant nominal (mode \emph{Heavy Duty}), la r\'ef\'erence retenue est l'\textbf{ATV630C11N4}, dont le \textbf{courant nominal est $I_n = \SI{170}{\ampere}$}.

Le variateur assure lui-m\^eme la protection thermique du moteur : le param\`etre $I_{th}$ (courant thermique) est r\'eglable de $0{,}2$ \`a $1{,}1\,I_n$, o\`u $I_n$ d\'esigne le courant nominal \textbf{du variateur}.

\begin{center}
\begin{tikzpicture}[>=Stealth, scale=0.85, every node/.style={scale=0.85}]
  \node[align=center,font=\scriptsize] (res) at (-0.2,0.6) {R\'eseau\\$\SI{400}{\volt}$};
  \node[draw,minimum width=2.0cm,minimum height=1.0cm,align=center,font=\scriptsize] (var) at (2.6,0.6) {ATV630C11N4\\$I_n{=}\SI{170}{\ampere}$\\\emph{prot. thermique}};
  \node[draw,minimum width=1.9cm,minimum height=1.0cm,align=center,font=\scriptsize] (mot) at (5.9,0.6) {MAS\\$P_u{=}\SI{90}{\kilo\watt}$\\$\eta{=}0{,}942$};
  \node[draw,circle,minimum size=0.7cm,font=\scriptsize] (p) at (8.2,0.6) {P};
  \draw[->] (0.8,0.6) -- (var.west);
  \draw[->] (var.east) -- (mot.west);
  \draw[->] (mot.east) -- (p.west) node[midway,above,font=\scriptsize]{pompe};
  \draw[dashed,->] (var.south) .. controls (4.2,-0.5) .. (mot.south) node[midway,below,font=\scriptsize]{$I_{th}$};
\end{tikzpicture}
\end{center}

\begin{enumerate}[leftmargin=1.5em]
  \item Le variateur est choisi sur le crit\`ere d'une surcharge de \SI{150}{\percent} du courant nominal du moteur. Expliquer \`a quelle contrainte de l'application (entra\^inement d'une pompe) correspond cette exigence, et pourquoi elle interdit de dimensionner le variateur au plus juste.
  \item Calculer le courant nominal $I_{n\,moteur}$ absorb\'e par le moteur \`a son r\'egime nominal.
  \item En d\'eduire la valeur de r\'eglage du param\`etre $I_{th}$ du variateur, exprim\'ee en fraction de $I_n$. V\'erifier qu'elle appartient bien \`a la plage r\'eglable et conclure sur le r\^ole de ce r\'eglage.
\end{enumerate}

}{
\textbf{Corrig\'e.}

\begin{enumerate}[leftmargin=1.5em]
  \item L'entra\^inement d'une pompe impose au variateur de d\'elivrer transitoirement un courant tr\`es sup\'erieur au courant nominal : au \textbf{d\'emarrage} (il faut vaincre l'inertie de la colonne d'eau et les frottements \`a l'amor\c{c}age) et lors des \textbf{surcouples} li\'es aux variations de charge hydraulique (variation de hauteur manom\'etrique, d\'emarrage sur conduite pleine). Le mode \emph{Heavy Duty} garantit une capacit\'e de surcharge de \SI{150}{\percent} du courant nominal moteur pendant une dur\'ee limit\'ee.

  Dimensionner le variateur au plus juste (sur le seul courant nominal permanent) conduirait \`a une \textbf{mise en s\'ecurit\'e du variateur \`a chaque d\'emarrage}, ou \`a un \'echauffement excessif de ses semi-conducteurs de puissance : la r\'eserve de surcharge fait partie int\'egrante du crit\`ere de choix.
  \item Le moteur \'etant triphas\'e, la puissance \textbf{absorb\'ee} vaut $P_a = P_u/\eta$ et s'exprime $P_a = \sqrt{3}\,U\,I\cos\varphi$ :
  \[ I_{n\,moteur} = \frac{P_u}{\eta\,\sqrt{3}\,U\cos\varphi} = \frac{90\,000}{0{,}942\times\sqrt{3}\times 400\times 0{,}86} \approx \SI{160,4}{\ampere} \]
  On v\'erifie la coh\'erence du choix du variateur : $I_n = \SI{170}{\ampere} > \SI{160,4}{\ampere}$, et la r\'eserve de surcharge de \SI{150}{\percent} porte la capacit\'e transitoire \`a environ \SI{255}{\ampere}, sup\'erieure aux $1{,}5\times 160{,}4 = \SI{241}{\ampere}$ demand\'es par le cahier des charges.
  \item Le param\`etre $I_{th}$ doit \^etre r\'egl\'e sur le courant nominal du \textbf{moteur} \`a prot\'eger, mais il s'exprime en fraction du courant nominal du \textbf{variateur} :
  \[ I_{th} = I_{n\,moteur} = k\,I_n \quad\text{avec}\quad k = \frac{I_{n\,moteur}}{I_n} = \frac{160{,}4}{170} \approx 0{,}94 \qquad\text{soit}\qquad \boldsymbol{I_{th} = 0{,}94\,I_n} \]
  Cette valeur appartient bien \`a la plage r\'eglable $[0{,}2\,I_n\,;\,1{,}1\,I_n]$, et se situe confortablement \`a l'int\'erieur de celle-ci.

  \smallskip
  \textbf{Conclusion.} Le r\'eglage $I_{th} = 0{,}94\,I_n$ cale exactement l'image thermique interne du variateur sur le courant nominal r\'eel du moteur (\SI{160,4}{\ampere}) et non sur celui du variateur (\SI{170}{\ampere}) : le moteur est ainsi prot\'eg\'e contre les surcharges prolong\'ees \emph{par le variateur lui-m\^eme}, ce qui justifie que le disjoncteur plac\'e en amont soit du type \textbf{MA (d\'eclencheur magn\'etique seul)}, sans relais thermique redondant. Laisser $I_{th}$ au maximum de $1{,}1\,I_n = \SI{187}{\ampere}$ reviendrait \`a autoriser en permanence \SI{17}{\percent} de surcharge sur le moteur sans aucune d\'etection : le moteur ne serait alors prot\'eg\'e par rien.

  \smallskip
  \textit{Point de vigilance jury (rapport de jury CAPET 2021, bilan de la partie C) : \guillemotleft\ Les candidats se sont limit\'es \`a proposer une r\'ef\'erence de composants parfois incompl\`ete sans justification, sans indiquer clairement les crit\`eres de choix. \guillemotright}
\end{enumerate}
}
\end{tcolorbox}

\CAPETENDFILE
